\section{Introduction}\label{sec:introduction}

Campana's theory of special varieties relates the absence of fibrations
of general type to restrictions on topology and monodromy
\cite{Campana2004,Campana2011}. Throughout the monodromy statements,
\emph{special} means special in the open-variety convention of
\cite[Definition 2.1]{CDY2025}: after any proper birational modification,
every dominant algebraic fibration with connected general fibre onto a
positive-dimensional normal quasi-projective base has orbifold-base
Kodaira dimension smaller
than the dimension of that base. The orbifold-base Kodaira dimension is
computed birationally,
using infimum multiplicities on smooth compactifications; a removed
divisor has multiplicity infinity. Section~\ref{sec:definitions}
spells out these conventions and the upper-model condition on orbifold
pairs that we verify for our surface example. That pair condition
implies the open-variety condition used here. For a smooth curve
$C=\bar C\setminus B$, specialness is equivalent to
\[
 \deg(K_{\bar C}+B)=2g(\bar C)-2+\#B\leq0.
\]
Thus elliptic curves are special, whereas a projective line with
at least three punctures is not.

In the compact setting, Campana's abelianity conjecture asks for a
virtually abelian ordinary fundamental group
\cite[Conjecture 7.1, p.~595]{Campana2004}. Campana--Claudon proved this
for special compact K\"ahler threefolds
\cite[Theorem 1.3]{CampanaClaudon2014}; the companion
\cite[Theorem 1.1]{OpenAICompactAbelianity2026} proves it in arbitrary
complex dimension. The open setting admits nonabelian nilpotent groups.
Our first result concerns complex linear images of ordinary fundamental
groups in this setting.

For a group $G$, put $\gamma_1G=G$ and
$\gamma_{r+1}G=[G,\gamma_rG]$. It has \emph{nilpotency class at most two}
if $\gamma_3G=1$, equivalently if every commutator is central. A property
holds \emph{virtually} if it holds in a subgroup of finite index.

\begin{theorem}[Two-step linear monodromy]\label{thm:linear}
Let $V$ be a connected smooth special complex quasi-projective variety.
For every positive integer $d$ and every representation
$\rho:\pi_1(V)\to\GL_d(\C)$, the image $\rho(\pi_1(V))$ has a subgroup
of finite index and nilpotency class at most two.
\end{theorem}

Cadorel--Deng--Yamanoi established virtual nilpotence of these linear
images: the identity component of their Zariski closure is a direct
product of a unipotent group and a torus
\cite[Theorem A/4.1]{CDY2025}. Cao--Deng--Hacon--P\u aun subsequently
announced the sharp class-two conclusion of Theorem~\ref{thm:linear}
\cite[Theorem A, equivalently Theorem 6.2]{CDHP2026}. Priority for that
conclusion belongs to their work. We give a different proof from the
independent CDY inputs and a structural theorem about unipotent
completion.

The bound is optimal. Cadorel--Deng--Yamanoi construct a smooth special
quasi-projective surface with a linear, two-step nilpotent fundamental
group that is not virtually abelian
\cite[Example 4.25 and Remarks 4.28--4.29]{CDY2025}. It is the complement
of the zero section in a degree-one line bundle on an elliptic curve;
its group is a Heisenberg central extension. Its ruled compactification
with the zero and infinity sections as reduced boundary has $K+D=0$
\cite[Lemma 5.10]{CDY2025}, and is special also in the orbifold-pair
convention by \cite[arXiv version, Theorem 6.7]{Campana2011}.

\subsection{The structural completion theorem}
A \emph{semiabelian variety} is a connected algebraic group fitting into
an exact sequence
\[
 1\longrightarrow\Gm^r\longrightarrow A\longrightarrow B
 \longrightarrow1,
\]
where $B$ is an abelian variety. For a connected smooth quasi-projective
variety $V$ with a chosen base point, the algebraic
\emph{quasi-Albanese map} $a_V:V\to\Alb(V)$ is the universal based
algebraic morphism to a semiabelian variety. This is the open analogue
of the Albanese map. We use the construction and universal property
recalled by Fujino \cite[Definition 2.18 and Theorem 3.16]{Fujino2024};
universality concerns algebraic morphisms, not arbitrary holomorphic
maps.

The \emph{rational unipotent completion} of a group $\Gamma$ is the
universal pro-unipotent algebraic group receiving a homomorphism from
$\Gamma$. Here pro-unipotent means an inverse limit of unipotent
algebraic groups, and universality means that every homomorphism from
$\Gamma$ to a finite-dimensional unipotent algebraic group over $\Q$
factors uniquely through the completion. For the finitely generated
groups considered here and each integer $c\geq1$, its
\emph{canonical $c$-step quotient} is the universal unipotent
quotient of nilpotency class at most $c$: every homomorphism to a
unipotent algebraic group of class at most $c$ factors uniquely through
it. Section~\ref{sec:mhs} gives its Lie-algebra construction.

\begin{theorem}[Completion under quasi-Albanese exactness]\label{thm:completion}
Let $V$ be a connected smooth complex quasi-projective variety and let
$a:V\to A=\Alb(V)$ be its algebraic quasi-Albanese map. Suppose $a$ is
dominant, choose a smooth connected general fibre $F$ and a point
$v\in F$, and assume that the sequence of ordinary fundamental groups
\begin{equation}\label{eq:exact}
 \pi_1(F,v)\longrightarrow\Gamma:=\pi_1(V,v)
 \xrightarrow{a_*}\Lambda:=\pi_1(A,a(v))\longrightarrow1
\end{equation}
is exact. Then every canonical rational unipotent quotient of $\Gamma$
has nilpotency class at most two. The full rational unipotent completion
is a finite-dimensional unipotent algebraic group of class at most two.
Every finite-dimensional complex unipotent representation of $\Gamma$
also has image of class at most two.
\end{theorem}

Exactness in \eqref{eq:exact} identifies the image of the fibre group
with the whole kernel of $a_*$; the first arrow need not be injective.
The theorem assumes neither specialness of $V$ or $F$ nor initial
nilpotence of the discrete group $\Gamma$. Its geometric hypotheses
are useful because the quasi-Albanese map identifies first homology,
and the algebraic fibre inclusion gives a morphism of pointed mixed
Hodge structures. These two facts connect the actual fibre image to
the commutator of each canonical unipotent quotient.

The proof uses Deligne's mixed Hodge theory for smooth open varieties
\cite{Deligne1971}, Morgan's foundational work on nilpotent homotopy data
\cite{Morgan1978,Morgan1986}, and Hain's canonical pointed construction
\cite{Hain1987}. At each finite stage, let $\g$ be the canonical
nilpotent Lie algebra and $\kk$ the image of the fibre Lie algebra.
Exactness and the quasi-Albanese first-homology isomorphism imply
$\kk=[\g,\g]$. The abelianization of $\kk$ is a quotient of
$H_1(F,\Q)$, with weights $-1$ and $-2$, while its position in a
commutator forces weights at most $-2$. It is therefore pure of weight
$-2$. The conjugation bracket has source of weights at most $-3$, so it
vanishes in this abelianization. This yields
$\gamma_3\g\subseteq\gamma_4\g$, and nilpotence forces both terms to
vanish. We first prove the bound on canonical quotients carrying mixed
Hodge structures, then pass to arbitrary unipotent representations by
universality.

To deduce Theorem~\ref{thm:linear}, we apply this structural result after
a finite \'etale cover. Such covers preserve specialness, and the
quasi-Albanese dominance and ordinary exactness results of
Cadorel--Deng--Yamanoi apply to the cover itself
\cite[Theorem 3.1, Lemma 4.4 and Proposition 4.13]{CDY2025}.
The resulting unipotent image has class at most two, and the torus
factor in the Zariski closure is abelian. This also explains why no
specialness assertion about the general fibre is needed.

There are several predecessors for the relation between Albanese maps
and nilpotent quotients. In the compact K\"ahler case, Campana proved
that surjectivity of the Albanese map makes every torsion-free
nilpotent quotient abelian \cite[Corollary 3.1]{Campana1995}.
His argument uses a homological criterion of Stallings. Aguilar
Aguilar--Campana obtained the same abelianity conclusion for smooth
quasi-projective varieties whose quasi-Albanese map is proper and
surjective \cite[arXiv version, Corollary 2(2)]
{AguilarAguilarCampana2025}. Theorem~\ref{thm:completion} permits a
nonproper map; the weight-$-2$ contribution from the fibre boundary
accounts for its class-two conclusion. When $H_1(F,\Q)$ is pure of
weight $-1$, the proof instead gives an abelian completion
(Remark~\ref{rem:proper-fibre}).

Rogov obtained another completion theorem through higher Albanese
manifolds \cite[Theorem B, equivalently Theorem 7.9 and Corollary 7.10]
{RogovHigher2026}. For a normal quasi-projective variety and a level
$s\geq3$, he assumes either dominance of its $s$th higher Albanese map
or a definable biholomorphism between that higher Albanese manifold and
the definable analytification of a quasi-projective variety. He proves
stabilization at level two. Theorem~\ref{thm:completion} instead starts
from ordinary fundamental-group exactness for the classical
quasi-Albanese map. Ordinary dominance alone is not substituted for
either set of hypotheses.

\subsection{A special surface with nonspecial quasi-Albanese fibre}
Campana asked whether general Albanese fibres of compact special
varieties are special \cite[Question 5.4, p.~577]{Campana2004}.
Campana--Claudon answered this for smooth projective varieties
\cite[Theorem 2.4]{CampanaClaudon2016}. The following open surfaces
show that the analogous assertion fails with reduced boundary. The
construction retains exceptional affine lines after deleting the
strict transforms of horizontal elliptic curves.

\begin{theorem}\label{thm:example}
Let $E$ be a complex elliptic curve, let $q_1,\ldots,q_n$ be distinct points
of $\PP^1$, where $n\geq3$, and choose points $e_1,\ldots,e_n$ of $E$.
Let
\[
 b:Y=\Bl_{\{(e_i,q_i)\}_{i=1}^n}(E\times\PP^1)
     \longrightarrow E\times\PP^1.
\]
Write $D_i$ for the strict transform of $E\times\{q_i\}$ and put
$D=\sum_iD_i$ and $X=Y\setminus D$. Then:
\begin{enumerate}[label=\textup{(\roman*)},leftmargin=2em]
\item $(Y,D)$ is a smooth projective special pair in the convention of
      \cite[Definitions 2.1--2.9]{CDHP2026}, including every upper model
      in that definition.
\item The morphism $a=(\operatorname{pr}_E\circ b)|_X:X\to E$ is the
      algebraic quasi-Albanese map, after choosing origins. Every
      fibre is connected, and its general fibre is
      $F\simeq\PP^1\setminus\{q_1,\ldots,q_n\}$.
\item The induced map $a_*:\pi_1(X)\to\pi_1(E)$ is an isomorphism.
      Thus $\pi_1(X)\simeq\Z^2$, and the homomorphism
      $\pi_1(F)\to\pi_1(X)$ is trivial.
\end{enumerate}
\end{theorem}

The general fibre has logarithmic canonical degree $n-2>0$, and so is
not special. This contradicts the fibre-specialness assertion of
\cite[Theorem E, equivalently Theorem 7.15]{CDHP2026} in the fixed
version arXiv:2603.14539v2. The fibre's ordinary group is free of rank $n-1$,
but its image in the total group is trivial. Thus the example
identifies the difference between the fibre and the image controlled
by Theorem~\ref{thm:completion}; the total group $\pi_1(X)=\Z^2$
satisfies the monodromy bound.

The geometric mechanism also explains the relevant orbifold
calculation. Above each $q_i$, the complete fibre has one deleted
component and one retained component, both of multiplicity one. The
retained component makes the direct orbifold multiplicity equal to
one. Contracting it before calculating the base forgets its
contribution. Campana's composition rule records precisely the
exceptional-divisor hypotheses needed for equality of these two bases
\cite[arXiv version, Proposition 3.13]{Campana2011}.
Section~\ref{sec:composition} identifies the resulting failure in
\cite[Claim 7.17]{CDHP2026}, where flatness over the Iitaka base is
used as flatness over a different, contracted intermediate space.

A related construction of Cadorel--Deng--Yamanoi also retains an
exceptional curve after deleting a boundary strict transform
\cite[Example 4.34]{CDY2025}. Here the product $E\times\PP^1$ produces
a positive-dimensional nonspecial quasi-Albanese fibre. Complete
elliptic slices and orbifold Riemann--Hurwitz verify specialness on
every required upper model. A complete section and the absence of
nonconstant units establish the quasi-Albanese universal property.
Local discs through the retained exceptional curves kill every
puncture loop and give the ordinary fundamental group.

Section~\ref{sec:mhs} proves Theorem~\ref{thm:completion}, and
Section~\ref{sec:linear} derives Theorem~\ref{thm:linear} and its
logarithmic-Kodaira-zero consequence.
Section~\ref{sec:definitions} sets out the orbifold conventions and
surface geometry. Sections~\ref{sec:specialness} and~\ref{sec:albanese}
prove Theorem~\ref{thm:example}; Section~\ref{sec:composition} explains
the composition calculation. Section~\ref{sec:ordinary} records the
remaining ordinary-group question. The representation results do not
assume that the ordinary group embeds in its completion, and no such
embedding is supplied by the argument.
